\section{Cross-Study Comparison and Evidence Gaps}
\label{sec:cross-study}

This section distills what the current evidence says across attack mechanisms. We use the core attack studies as the denominator for cross-study claims and use contextual references only to explain substrates and evaluation concepts. The key point is that attacks should be compared by mechanism, validation tier, action representation, and consequence, not by a single undifferentiated ASR.

\begin{table*}[t]
\centering
\scriptsize
\caption{Descriptive statistics from frozen database \texttt{VLA-Attack-ER-v2026-06-16}. Attack-study statistics use the Core Evidence Corpus unless a direct-VLA subset is explicitly named; NR is counted as not reported.}
\label{tab:evidence-statistics}
\renewcommand{\arraystretch}{1.08}
\setlength{\tabcolsep}{3pt}
\begin{tabularx}{\textwidth}{@{}p{0.22\textwidth}Y@{}}
\toprule
\textbf{Statistic} & \textbf{Observed pattern} \\
\midrule
Corpus split & 85 records: 30 core attack studies (23 direct VLA and 7 embodied-adjacent) plus 55 contextual references. \\
Direct VLA subset & 23/30 core attack studies are direct VLA evidence; all percentages below use the direct subset ($n=23$). \\
Entry surface flags in direct VLA & Visual/physical visual input: 14/23 (60.9\%); language/instruction: 9/23 (39.1\%); reasoning-trace input: 1/23 (4.3\%); privacy/query output: 2/23 (8.7\%); poisoned trajectory/data trigger: 1/23 (4.3\%). Multi-label counts overlap. \\
Knowledge level & Any black-box or transfer-only condition: 13/23 (56.5\%); any white/gray-box optimization condition: 14/23 (60.9\%). Multi-label counts overlap when papers report multiple settings. \\
Access and privilege & Digital access channel: 14/23 (60.9\%); physical or digital-to-physical channel: 8/23 (34.8\%); supply-chain or data-contributor channel: 5/23 (21.7\%). \\
Highest validation tier & Digital input-only: 3/23 (13.0\%); simulation rollout: 10/23 (43.5\%); real-sensor input: 1/23 (4.3\%); real-robot actuation: 9/23 (39.1\%); HIL: 0/23 (0\%); human-proximate: 0/23 (0\%). \\
Model/benchmark concentration & OpenVLA-family target appears in 10/23 (43.5\%); LIBERO appears in 12/23 (52.2\%); $\pi_0/\pi_0$-fast appears in 6/23 (26.1\%). \\
Reporting practices & Explicit transfer beyond one target/model/prompt: 16/23 (69.6\%); limited transfer notes: 6/23 (26.1\%); benign utility reported: 7/23 (30.4\%); query/edit budget reported: 6/23 (26.1\%). \\
Robot-control reporting & Action/chunk horizon reported: 3/23 (13.0\%); controller frequency reported: 1/23 (4.3\%); safety filter reported: 0/23 (0\%); recovery policy/outcome reported: 0/23 (0\%); pre/post-controller trace: 0/23 (0\%). \\
\bottomrule
\end{tabularx}
\end{table*}


\begin{table*}[t]
\centering
\scriptsize
\caption{Representative direct VLA attack mechanisms and what each mechanism currently supports. The table is qualitative: it summarizes mechanism-level evidence rather than pooling heterogeneous ASR values.}
\label{tab:core-heatmaps}
\renewcommand{\arraystretch}{1.08}
\setlength{\tabcolsep}{3pt}
\begin{tabularx}{\textwidth}{@{}p{0.17\textwidth}Y Y Y Y@{}}
\toprule
\textbf{Attack family} & \textbf{Representative entry} & \textbf{Changed representation} & \textbf{Robotic consequence} & \textbf{Evidence strength and caveat} \\
\midrule
Instruction/control & Prompt edit, scene text, or cross-modal instruction perturbation & Plan, source grounding, action token, or command-preserving trajectory & Task failure, action inflation, wrong-object behavior, or trajectory redirection & Direct rollout evidence is emerging; transfer, query budget, and recovery still vary across studies. \\
Visual/physical & Digital patch, printed patch, object texture, or partial-observation physical access & Embedding, grounding, visual proprioception, action loss, or trajectory objective & Wrong-object action, targeted motion, trajectory deviation, or degraded manipulation & Strong mechanism evidence across digital, simulation, and some real-robot settings; physical generality across viewpoints, objects, and controllers remains open. \\
Backdoor/action & Poisoned demonstration, trigger object, or supply-chain fine-tuning & Primitive, action freeze, chunk drift, waypoint, or long-horizon action sequence & Triggered release/freeze, wrong goal behavior, or gradual trajectory shift & Direct VLA studies show action representations are attackable while preserving clean utility; robustness to filtering, adaptation, and controller clipping is less mature. \\
Reasoning/privacy & Physical/contextual reasoning perturbation or action-trace query access & CoT/entity reference, reasoning-to-action state, likelihood, attention, or temporal trajectory signature & Reasoning hijack, unsafe driving/manipulation behavior, or training-membership leakage & Expands the field beyond task-failure ASR; privacy and reasoning claims require metrics different from physical task success. \\
\bottomrule
\end{tabularx}
\end{table*}


\subsection{What Current Evidence Tells Us}

The first lesson is that VLA attacks are increasingly \emph{representation attacks}, not only input perturbations. Visual patches matter because they can move embeddings, grounding, visual proprioception, and trajectory objectives; instruction attacks matter because text can steer plans and action tokens; backdoors matter because triggers can bind to reusable action primitives or chunks. This explains why a survey organized only by input modality misses the robotics mechanism.

The second lesson is that closed-loop evidence is still thin. Many papers show that an attack changes a rollout, but fewer show what happens between the VLA output and the physical state: controller clipping, safety filtering, latency, emergency stop, or recovery. Current evidence therefore supports claims about vulnerable VLA representations more strongly than claims about deployed robot hazard rates.

The third lesson is that the field is splitting into two complementary directions. One direction seeks stronger and more realistic attackers: black-box instruction edits, transferable patches, object-bound 3D textures, partial-observation attacks, and supply-chain triggers. The other direction seeks better evaluation: typed ASR, clean utility, transfer axes, validation tiers, and privacy metrics. A mature VLA attack literature will need both; stronger attacks without robotics-aware evaluation will not answer T-RO-relevant safety questions, and evaluation protocols without adaptive attacks will underestimate residual risk.

\subsection{Direct VLA Evidence by Attack Mechanism}

The direct matrix separates four mechanism families that are often conflated. First, instruction-control attacks operate through language while measuring robot behavior: RoboGCG asks whether adversarial text can give low-level control authority \cite{jones2025robogcg}; SABER asks whether bounded black-box edits can induce task failure, action inflation, or constraint violations across VLA policies \cite{wu2026saber}; Q-DIG asks whether natural, task-relevant prompt diversity can systematically expose VLA policy failures \cite{srikanth2026qdig}; and trajectory redirection asks whether a command-preserving prompt can move the full rollout toward an attacker-specified physical outcome \cite{puthumanaillam2026trajectory}. These are all language-entry attacks, but they differ in objective, attacker cost, and whether the attack is targeted.

Second, visual attacks split into action-loss, representation-loss, transfer, and physical-object variants. Wang et al. optimize VLA-specific action discrepancy and trajectory objectives \cite{wang2025vlaadversarial}. EDPA and ADVLA move the target earlier, toward visual embeddings, projected features, attention, and vision--language alignment \cite{xu2025edpa,zhang2025advla}. UPA-RFAS and VLA-Hijack study transfer by exploiting shared feature-space, attention, semantic, or visual-proprioception structure rather than only a model-specific action head \cite{lu2026patch,fu2026vlahijack}. Tex3D and partially observable patch attacks add physical realism dimensions: object-bound textures, trajectory-aware optimization, and prefix-limited attacker observation \cite{chen2026tex3d,wang2026partialpatch}.

Third, action/backdoor attacks are not merely training-time versions of patches. BadVLA targets trigger-conditioned action deviation \cite{zhou2025badvla}; GoBA targets physical-object-triggered goal behavior \cite{zhou2025goba}; DropVLA targets a reusable primitive at a short reaction window \cite{xu2025dropvla}; AttackVLA compares multiple attack implementations and introduces long-horizon triggered action sequences \cite{li2025attackvla}; FreezeVLA and SilentDrift target action freezing and chunk-level drift \cite{wang2025freezevla,xu2026silentdrift}. The important comparison variable is the attacked action representation: action token, gripper primitive, static action state, full trajectory, or chunked delta-pose sequence.

Fourth, reasoning and privacy attacks expand the objective set. TRAP and altered-thought evaluations show that reasoning traces, entity references, and CoT-to-action pathways can be attack targets in reasoning-enabled VLA systems \cite{huang2026trap,altered2026thoughts}. ReasonBreak shows a related reasoning--trajectory vulnerability in autonomous-driving VLA systems, where closed-loop collision and safety metrics matter more than tabletop task success \cite{teymoorianfard2026reasonbreak}. Membership inference and VLALeaks show that action likelihoods, action errors, temporal motion patterns, attention discrepancies, and trajectory structure can leak training membership \cite{peng2026miavla,luan2026vlaleaks}. These attacks should be evaluated with confidentiality metrics, not forced into task-failure ASR.

\subsection{What Can and Cannot Be Compared}

ASR is not directly comparable across the current VLA attack literature. In one paper, ASR may mean task failure under a visual patch; in another, targeted gripper release within a reaction window; in another, action inflation; in another, trajectory redirection; and in privacy work, AUC or TPR@FPR for membership inference. The denominator also varies: image frames, task episodes, triggered episodes, candidate trajectories, model queries, or physical trials. Validation tier varies from digital input-only to simulation rollout to real-robot actuation. Pooling these numbers would create false precision.

Comparable claims are therefore qualitative and typed. Direct evidence supports the claim that VLA policies can be attacked through vision, language, action representations, backdoors, reasoning traces, and privacy channels. It does not yet support a universal claim that any particular attack transfers across all VLA architectures, embodiments, controllers, or safety envelopes. Embodied-adjacent evidence supports risk hypotheses about source confusion, tool misuse, middleware compromise, sensor corruption, and HRI authority ambiguity. Analogical evidence supports mechanisms and evaluation requirements, not VLA vulnerability claims by itself.

\subsection{Evidence Gaps}

The strongest direct evidence is concentrated in OpenVLA and LIBERO-style manipulation. This is useful because it creates a common baseline, but it weakens claims about general VLA security. Flow/diffusion policies, bimanual systems, mobile manipulation, navigation, tactile policies, and high-level tool/skill interfaces are less consistently covered. Action representations are also unevenly studied: tokenized action heads and patch-induced trajectory deviation are common; diffusion/flow denoising paths, velocity fields, controller latency, temporal ensembling, recovery policies, and safety-envelope interactions are much less mature.

Physical validation is improving but remains narrow. Some direct VLA papers include physical robot demonstrations, but few provide enough detail on sensor stack, control frequency, calibration, lighting, object mass, latency, recovery behavior, and safety filters to make deployment claims portable. Hardware-in-the-loop is surprisingly underdeveloped as a middle tier between simulation and full physical trials.

Privacy is an emerging direct VLA attack area rather than a side issue. Membership inference papers show that VLA training data can leak through action outputs, likelihoods, attention patterns, and temporal motion signatures \cite{peng2026miavla,luan2026vlaleaks}. This evidence is direct, but it supports confidentiality conclusions, not physical-safety conclusions. Future survey tables should therefore include privacy objectives rather than treating all VLA attacks as task-failure attacks.

\subsection{Target Model and Benchmark Concentration}

The target-model distribution is informative evidence. OpenVLA and OpenVLA-OFT appear frequently because they are open, strong, and easy to evaluate on LIBERO \cite{kim24openvla,kim2025openvlaoft,liu2023libero}. $\pi_0$ and $\pi_0$-fast appear in several attacks because flow or diffusion-style action generation changes both inference speed and attack target representation \cite{black2024pi0}. SpatialVLA appears in benchmark-oriented attack work because spatial action grids and 3D spatial encoding create a different target representation from plain tokenized actions \cite{qu2025spatialvla}. UniVLA, CronusVLA, reasoning VLAs, Alpamayo-style driving VLAs, and VLA-Adapter broaden the evidence base, but they are not yet as repeatedly tested as OpenVLA-style manipulation systems \cite{fu2026vlahijack,huang2026trap,teymoorianfard2026reasonbreak,xu2026silentdrift}.

The benchmark distribution is similarly concentrated. LIBERO is the dominant manipulation benchmark in the direct attack matrix, with RLBench, CALVIN, ManiSkill, RoboCasa, and DROID-derived data serving more often as capability or transfer context \cite{james2020rlbench,mees2022calvin,mu2021maniskill,nasiriany2024robocasa,khazatsky2024droid}. This concentration is useful for comparability but risky for authority: a benchmark can overrepresent tabletop manipulation, fixed cameras, short horizons, and simplified safety envelopes. Benchmark audits and robot-data papers therefore remain relevant to attack evidence even when they are not attacks themselves \cite{jiang2026benchmarkaudit,openx2023,walke2023bridgedata,dasari2019robonet}.

\subsection{Reporting Coverage}

The direct VLA literature reports clean utility more often for backdoors than for inference-time attacks, because a backdoor must remain hidden under benign inputs. Transfer is reported most often in visual patch and benchmark papers, but transfer axes differ: model transfer, task transfer, viewpoint transfer, suite transfer, and sim-to-real transfer are not interchangeable. Query budget appears in black-box instruction work but is absent from many visual and backdoor studies. Adaptive settings remain rare: most papers test a fixed attack against an undefended or lightly defended policy, not against a deployed prompt hierarchy, source-authentication layer, action shield, or runtime monitor. Recovery is the weakest reporting dimension. Even when physical trials are included, many papers end evaluation at task failure or target action induction rather than measuring whether the robot could stop, replan, ask for help, or recover under controller constraints.

\subsection{Open Versus Understudied}

Some topics are genuinely open: multi-session memory poisoning of VLA robots, adaptive attacks against deployed safety envelopes, HRI attacks with ethically controlled human subjects, and standardized HIL protocols. Other topics have preliminary direct evidence but insufficient breadth: transferable visual patches, black-box instruction attacks, action-level backdoors, reasoning-chain attacks, and VLA privacy attacks. These should be described as early but real VLA evidence, not as unexplored.
